\section{Introduction}\label{sec:r-introduction}

A finite message chosen from data changes the conditional distribution
of those data.  This elementary fact becomes a central issue in
memory-constrained inference.  Suppose a signal \(S\) is observed
through exact Gaussian equations and a learner retains a finite state
\(W\).  After conditioning on \(W\), the rows used to select that
state are generally biased.  An independent Gaussian matrix has the
same unconditional row law, but it need not leave the same information
about \(S\) once its labels are revealed.

This paper quantifies that change.  We compare information beyond a
wholly independent auxiliary projection with information beyond a
projection containing the actual rows used by the learner.  The two
experiments keep the signal--message marginal and the signal--matrix
marginal fixed.  Their full joint laws differ.  At suitable row
dimensions proportional to \(d\), the comparison cost is \(O(d)\)
when the message entropy is at most \(d^2\).  We also prove the
projection moment that makes this possible and develop two further
comparisons based on replacing one row at a time and on a
two-signal fiber measure.

Here is one precise comparison.  Write \(H\) for discrete entropy
in nats, \(h\) for differential entropy, and
\(I(\cdot;\cdot\mid\cdot)\) for conditional mutual
information.  Let \(S\) be uniform on \(S^{d-1}\), put
\(m=\lfloor d/10\rfloor\) and \(\ell=\lfloor d/2\rfloor\), and
let \(A\) be an \(m\)-row standard Gaussian matrix independent of
\(S\).  A countable message \(W\) may have arbitrary joint
dependence on \((S,A)\), subject to \(H(W)\le d^2\).
Draw an \((\ell-m)\)-row Gaussian matrix \(C\) independently of
\((S,A,W)\), and an \(\ell\)-row Gaussian matrix \(G\) independently of
\((S,W)\).  Theorem~\ref{thm:r-critical-comparison} proves
\[
 I(S;W\mid G,GS)
 \le I(S;W\mid A,C,AS,CS)+Kd
\]
for all sufficiently large \(d\), with an absolute threshold and
an absolute constant \(K\).  Thus replacing rows that helped select
the message by fresh rows costs only \(O(d)\) additional conditional
information, despite the possible bias of \(A\) given \(W\).

\subsection*{Memory, exact equations, and prior methods}

Memory restrictions affect statistical estimation even when
computation is unrestricted.  Steinhardt and Duchi established
memory-dependent minimax bounds for sparse noisy regression
\cite{SteinhardtDuchi2015}.  Raz's branching-program lower bounds
for parity learning and broader discrete learning tasks showed how
finite width can force a large sample cost \cite{Raz2016,Raz2017}.
Sharan, Sidford, and Valiant developed a closely related finite-state
framework for continuous regression \cite{SSV2019}.

Their precision theorem uses isotropic Gaussian covariates with
independent uniform additive noise of half-width \(2^{-d/5}\).
For Euclidean accuracy \(d^{-r}\), with \(r\le O(d/\log d)\),
at most \(d^2/4\) bits of memory, and success probability at least
\(2/3\), it gives an \(\Omega(d\log r)\) sample lower bound
\cite[Theorem~1]{SSV2019}.  The present observation law has exact
labels, which contain more information than noisy labels.
Accordingly, our proof needs an estimate for exact projection
densities.  Sharan, Sidford, and Valiant's expansion of a projection
moment into independent signal copies is a direct method precedent
\cite[Section~7, Lemma~12]{SSV2019}; their one-row interval estimate
uses a positive noise scale and global \(L^2\) control, while the
mixed multirow estimate below is proved from local mass bounds.
Dagan, Kur, and Shamir give related space lower bounds for an
approximate-nullvector task with independent Gaussian rows and for
empirical regression on a randomly ordered finite data set with
exactly consistent labels \cite[Theorems~3 and~10]{DKS2019}.
The latter result bounds the empirical residual under a hard joint
data distribution; our recovery problem uses independent Gaussian
rows and a fixed signal.

Information-theoretic analyses of adaptive selection also compare a
selected statistic with an independent reference.  Russo and Zou
bound the divergence between a selected statistic on the observed
data and its evaluation on an independent copy by the information
used for selection \cite[Section~2.4, Proposition~4]{RussoZou2016}.
Xu and Raginsky bound an expectation difference between a joint law
and the product of its marginals for a test that is sub-Gaussian
under the product law
\cite[Section~3.1, Lemma~1]{XuRaginsky2017}.  These comparisons
explain the role of an information charge for selection.  Here the
test is an exact projection density, and a locally proved \(q\)th
moment makes the row-information charge appear with coefficient
\(1/q\).

The geometric calculation also belongs to the tradition of
controlling projection energies by inverse distances.  Mattila's
two-point averaging estimate gives a classical instance of that
principle \cite[proof of Lemma~5.1, p.~236]{Mattila1975}.  Our mixed moment
contains both inverse lengths and inverse distances to the span of
earlier signal differences.  The specific Gaussian factors,
dimension conditions, and finite fiber normalization are proved in
this paper.

\subsection*{The streaming model and the application}

The comparison theorems concern finite messages with arbitrary
dependence on a signal and some Gaussian rows.  Their streaming
application uses the following model.

\begin{definition}[Finite persistent-state learner]
\label{def:r-streaming-model}
Fix \(d\ge2\), an integer memory bound \(M\ge0\), an angular
accuracy \(0<\epsilon\le1/10\), and a deterministic finite integer
horizon \(T\ge0\).  The signal is a unit vector.  It may be fixed
before sampling, or sampled from a specified prior; the
average-case statements in this paper use uniform probability on
\(S^{d-1}\).  Pre-generate independent
\[
 X_t\sim N(0,I_d),\qquad Y_t=\langle X_t,S\rangle,\qquad t\ge1.
\]
The row sequence is independent of \(S\).

Let \(\Xi\) be a shared rule seed, constant if no shared seed is
used, and let \(U_0\) be the initial finite state.  The pair
\((\Xi,U_0)\) is jointly independent of
\((S,(X_t)_{t\ge1})\); \(\Xi\) and \(U_0\) may depend on one
another.  The original randomized experiment is jointly measurable
in the signal, the complete row sequence, the shared seed,
initialization, and all transition and output randomness.  All
variables and kernels are taken on standard Borel spaces, with the
completed-measurability convention below for sample arguments.

At every index, including index zero, the persistent state has at most \(2^M\)
possible values.  Conditional on \(\Xi\), the rule at time \(t\)
uses only the current state, the index \(t\), and the entire current
pair \((X_t,Y_t)\), together with fresh randomness that has no
additional signal or sample information.  The rule may perform
unrestricted computation on that pair, but everything retained for
later samples must be encoded in the new finite state.  The learner
receives the rows once in their prescribed order and cannot choose a
row or revisit an earlier one.

The learner stops at an index \(\tau\in\{0,\ldots,T\}\).  Conditional
on the fixed seed, its unit-vector output \(\widehat S\) is drawn
from a kernel depending only on the terminal state and \(\tau\),
using fresh output randomness.  In particular the last observed pair
can affect the output only through that terminal state.  The seed
may select nonuniform rules depending on \(d,\epsilon,t\), but it
has no data-dependent content.
Before the first sample, a stopping choice may use only the seed,
the initial state, and fresh randomness.

For each fixed seed, state, and index, the transition and stopping
probabilities may be Borel functions of the sample pair, or
measurable in the completion of Gaussian row measure times Lebesgue
label measure.  This per-seed convention supplements the required
joint measurability of the original mixed experiment; it does not
replace it.  Success means
\(\arccos\langle\widehat S,S\rangle\le\epsilon\).
\end{definition}

The joint condition on \((\Xi,U_0)\) ensures that fixing the rule
seed leaves initialization data-independent.  We record this
reduction because the information potentials start from the
conditional initial-state law.

\Needspace{6\baselineskip}
\begin{lemma}[Fixing the shared rule seed]\label{lem:r-seed-reduction}
In the model above, for almost every seed value \(\xi\),
\begin{equation}\label{eq:r-seed-factorization}
 \mathcal L\bigl(S,(X_t)_{t\ge1},U_0\mid\Xi=\xi\bigr)
 =\mathcal L\bigl(S,(X_t)_{t\ge1}\bigr)
       \otimes\mathcal L(U_0\mid\Xi=\xi).
\end{equation}
If the original uniform-prior success probability is at least
\(p_0\), one may choose such a seed value whose conditional
uniform-prior success is at least \(p_0\), with the same state
bound and horizon.  This assertion is about the uniform-prior
experiment.
\end{lemma}

\begin{proof}
Disintegrate the joint independence of
\((\Xi,U_0)\) and \((S,(X_t)_{t\ge1})\) over \(\Xi\).
It gives \eqref{eq:r-seed-factorization} for almost every seed.
Joint measurability of the original experiment makes its conditional
success probability a measurable function of \(\xi\), whose mean
is the original success probability.  Averaging therefore supplies
a seed with the stated success and factorization.  The conditional
initial-state law is retained.  A later Borel replacement is
performed separately in this fixed-seed experiment; no jointly
measurable selection of replacement versions is required.
\end{proof}

The critical-radius comparison gives the following streaming
consequence: for \(M(d)=o(d^2)\), every accuracy sequence
\(0<\epsilon(d)\le1/10\), and uniform-sphere average success at
least \(3/5\), the deterministic horizon must satisfy
\[
 T(d)=\Omega\!\left(d\log\frac1{\epsilon(d)}\right).
\]
The lower-bound constant is absolute; the eventual dimension
threshold may depend on the memory sequence.  Section~\ref{sec:r-streaming}
proves this statement and records the distinct constants and
thresholds arising from the other comparisons.  A learner that
succeeds with probability at least \(2/3\) for every fixed signal
also has that success under the uniform prior, so the worst-signal
consequence follows separately.  The auxiliary projections used in
the proof are analysis variables; the learner receives only the
exact stream in Definition~\ref{def:r-streaming-model}.

\subsection*{How the comparisons work}

Section~\ref{sec:r-experiments} defines the aligned and independent
experiments explicitly.  Their information difference contains a
negative conditional row-information term.  Data processing
compares that term with the divergence between the actual
conditional label law and the projection of the posterior
\(P_{S\mid W}\).  What remains is the entropy of the independent
projection plus an aligned log density.

Section~\ref{sec:r-mixed-moments} proves the analytic estimate for
that density.  In a \(q\)th moment, the actual rows are shared by
all \(q\) factors while independent auxiliary rows are averaged
separately in each factor.  Gaussian integration produces inverse
lengths and the inverse Gram determinant of the signal differences.
A local mass bound controls thin tubes around their successive
spans, and a second bound controls the sum over small radial scales.
The section treats exact density versions and the order of the
change of row law and limiting operation.

Sections~\ref{sec:r-critical-radius} and
\ref{sec:r-density-levels} use this estimate after refining a
posterior by its local mass.  The critical-radius refinement records
both a radius and its logarithmic mass level.  Two further
comparisons use a single maximal level or a pair of density labels.
Their entropy costs are explicit, and the density-level terms
cancel with opposite signs.  The two-label argument also retains
its own Gaussian-mollifier and first-density passage.

These comparisons have different scopes.  The critical-radius result
allows countable messages and \(\lfloor d/10\rfloor\) actual rows.
The two-label result treats finite messages with the larger actual
block \(\lfloor d/8\rfloor\); its Gaussian-mollifier density version
supports an estimate at every fixed target.  The one-level result
uses only one refinement label, a net-weight entropy estimate, and
\(\lfloor d/32\rfloor\) actual rows.  It keeps the entropy-dependent
comparison cost explicit rather than imposing \(H(W)\le d^2\)
in its statement.

The last two comparisons use different geometric objects.
Section~\ref{sec:r-hybrids} keeps the actual ending state fixed
while replacing an actual prefix of rows by an independent suffix,
one row at a time.  Independent random kernels and a normalized
reciprocal-distance selection control the successive changes.
Section~\ref{sec:r-fiber} constructs a finite measure on pairs
having the same Gaussian label.  It proves that measure's exact
kernel and normalizes the reference before applying probability
data processing.  Local-mass ratios then bound the resulting
information difference.
The hybrid argument averages over the number of auxiliary observations,
whereas the fiber comparison concerns a state transition through an
actual block, rather than an arbitrary message jointly distributed
with its rows.  These scopes explain the separate applications below.

Finally, Section~\ref{sec:r-streaming} applies these comparisons
to successive fresh blocks.  It includes a local Borel-version
reduction for the fixed-seed uniform-prior experiment, the state
and stopping-index costs, and the residual-sphere information
required by an accurate output.  A separate memory-unrestricted
sample bound absorbs the final-block rounding cost.  This separates the
analytic comparisons from the bookkeeping needed to apply them
to a learner.
